\scanpage{60}
$\operatorname{int}(\sigma_\infty),\operatorname{int}(\tau_\infty)$,
\inserted{Or dans ce cas $g_0=\tau_\infty$ resp. $g_0=\sigma_\infty$, \deleted{\unsure{donc}} \inserted{on} trouve $r'=(i,\omega,\varepsilon')$ resp. $r'=(i,\omega',\varepsilon')$, donc}
\[
(39)\qquad\begin{cases}
\psi_i^{\omega,\varepsilon'}=\tau\circ\psi_i^{\omega,\varepsilon}=\psi_i^{\omega,\varepsilon}\circ\tau_\infty\\[5pt]
\psi_i^{\omega',\varepsilon'}=\underbrace{\sigma_i}_{\substack{\text{transposition}\\\in\mathcal G_J\text{ qui}\\\text{fixe }i}}\circ\psi_i^{\omega,\varepsilon}=\psi_i^{\omega,\varepsilon}\circ\sigma_\infty
\end{cases}
\]
donc en composant
\[
(40)\qquad\psi_i^{\omega',\varepsilon\deleted{\unsure{'}}}=\deleted{\unsure{\tau}}\underbrace{\widetilde\sigma_i}_{\substack{\tau\sigma_i\\\text{symétrie}\\\text{autour du}\\\text{grand cercle}\\\text{passant par}\\P_\varepsilon,P_{\varepsilon'},Q_i}}\circ\psi_i^{\omega,\varepsilon}=\psi_i^{\omega,\varepsilon}\circ\widetilde\sigma_\infty
\]
Dans tous ces calculs, suivant sans doute d'anciennes habitudes, j'ai privilégié les $P_\varepsilon$ et les $Q_i$ comme p\textsuperscript{ts} base -- alors que pour le travail que j'ai à faire, ce sont les points $R_i^\omega$ qui sont les plus cruciaux, où les calculs faits en ces p\textsuperscript{ts} capturent directement les phénomènes de ``monodromie'' autour des ``trous'' $R_i$. Je vais réparer cette négligence, en introduisant, en
\margin{59}
